\clearpage\NoSourceHead\SourceLeaf{397}
\SourceGroup{II}
\HistoricTitle{GRUNDBEGRIFFE DER LEHRE VON DEN REIHEN.}
\HistoricCitation{[Handschrift von 6 Seiten in Fa, Kapsel 46 a.]}
\begin{center}1.\end{center}
\noindent\hspace*{\parindent}Der Inbegriff einer jeden Anzahl beliebiger Grössen kann im weitern Sinne des Worts \textls[120]{eine Reihe} genannt werden; indess würde diese Ausdehnung des Begriffs wenig Nutzen haben und man beschränkt daher in der höhern Mathematik den Ausdruck \textls[120]{Reihe} auf den Inbegriff solcher Grössen, die, insofern man jeder derselben ihre eigene \textls[120]{Stelle} anweiset, d. i. die erste, zweite, dritte Grösse u. s. f. unterscheidet. alle \textls[120]{nach einem Princip} bestimmt werden. Der wesentliche Charakter einer Reihe ist also, dass für jeden Ort in derselben die entsprechende Grösse (das \textls[120]{Glied} der Reihe) sofort völlig bestimmt ist, und es daher als möglich angesehen wird, sobald man das Princip nach welchem die Reihe gebildet wird (ihr \textls[120]{Gesetz}) kennt, sie soweit man will fortzusetzen. Den Ort eines jeden Gliedes bezeichnet man durch eine Zahl, die der \textls[120]{Index} desselben heisst, so dass 1 der Index des ersten Gliedes ist, 2 der des zweiten u. s. f.
\begin{center}2.\end{center}
\noindent\hspace*{\parindent}Der hier gegebene Begriff einer Reihe ist von weiterem Umfange als der gewöhnliche, da man sie durch den Inbegriff der Werthe einer Function Einer veränderlichen Grösse erklärt, welche dieselbe erhält, indem man diese veränderliche Grösse nach und nach \(=1,\,2,\,3\) u s. f. setzt, wenn man anders nicht auch den Ausdruck Function in einer ausgedehnteren Bedeutung nehmen
\clearpage\SourceLeaf{398}\ArithmeticHead{391}{R}{GRUNDBEGRIFFE DER LEHRE VON DEN REIHEN.}
\noindent will, als bisher geschehen ist. Denn nach dieser Erklärung wäre z. B. die Reihe der Primzahlen \(1,\,2,\,3,\,5,\,7\) u. s. f. oder jede daraus abgeleitete wie
\[
1,\ \frac14,\ \frac19,\ \frac1{25},\ \frac1{49},\ \frac1{121}\ \text{u. s. f.}
\]
\noindent nicht unter den Reihen enthalten, welches dem Sprachgebrauche weniger gemäss ist. Zum Unterschiede kann man solche Reihen, wo jedes Glied durch eine analytische Function des Index dargestellt wird, \textls[120]{analytische} Reihen nennen, so wie man die Function selbst, allgemein vorgestellt. das \textls[120]{allgemeine Glied} der Reihe zu nennen pflegt.

